% from ega4-2-fr.tex
\oldpage[IV-2]{103}

tel idéal est nécessairement de hauteur 1; en outre, comme $E$ est une
$C$-algèbre finie et que $\mathfrak p\not\supset\mathfrak n$,
$C_{\mathfrak p}$ est intégralement clos (Bourbaki, \emph{Alg. comm.},
chap. V, § 1, n° 5, cor. 5 de la prop. 16), donc
$E_{\mathfrak p'}=C_{\mathfrak p}$, ce qui prouve nos assertions.

Il serait intéressant de savoir si tout anneau local intègre noethérien
vérifiant la conclusion de (5.6.10) est universellement caténaire; c'est ce
qu'a affirmé Nagata [33], mais sa démonstration ne paraît pas complète.
\end{example}

\subsection{\texorpdfstring{Profondeur et propriété $(S_k)$}{Profondeur et propriété (S\string_k)}}

\begin{definition}[5.7.1]
\label{IV.2.s5.7.1-fr}
Soient $X$ un préschéma localement noethérien, $\mathscr F$ un
$\mathcal O_X$-Module cohérent. On appelle \emph{profondeur} (resp.
\emph{coprofondeur}) de $\mathscr F$ en un point $x\in X$ le nombre
$\operatorname{prof}(\mathscr F_x)$ (resp.
$\operatorname{coprof}(\mathscr F_x)$) (0, 16.4.5 et 16.4.9). On appelle
\emph{coprofondeur de $\mathscr F$} le nombre
\begin{equation}
\operatorname{coprof}(\mathscr F)
=\sup_{x\in X}\operatorname{coprof}(\mathscr F_x).
\tag{5.7.1.1}\label{IV.2.s5.7.1.1-fr}
\end{equation}

On dit que $\mathscr F$ est un $\mathcal O_X$-Module de Cohen-Macaulay au
point $x$ si $\mathscr F_x$ est un $\mathcal O_x$-module de Cohen-Macaulay,
c'est-à-dire (0, 16.5.1) si $\operatorname{coprof}(\mathscr F_x)=0$. On dit
que $\mathscr F$ est un $\mathcal O_X$-Module de Cohen-Macaulay s'il l'est
en tout point, autrement dit si $\operatorname{coprof}(\mathscr F)=0$. Un
point $x\in X$ tel que $\mathcal O_x$ soit un anneau de Cohen-Macaulay est
encore appelé \emph{point de Cohen-Macaulay} de $X$.

On appelle \emph{coprofondeur de $X$} et on note $\operatorname{coprof}(X)$
le nombre $\operatorname{coprof}(\mathcal O_X)$. On dit que $X$ est un
\emph{préschéma de Cohen-Macaulay} si $\mathcal O_X$ est un
$\mathcal O_X$-Module de Cohen-Macaulay, autrement dit si
$\operatorname{coprof}(X)=0$. Tout préschéma localement noethérien de
dimension 0 est évidemment un préschéma de Cohen-Macaulay. Dire que
$\operatorname{Spec}(A)$ est un schéma de Cohen-Macaulay signifie que $A$
est un anneau de Cohen-Macaulay (0, 16.5.13).
\end{definition}

\begin{definition}[5.7.2]
\label{IV.2.s5.7.2-fr}
Soient $X$ un préschéma localement noethérien, $\mathscr F$ un
$\mathcal O_X$-Module cohérent, $k$ un entier positif ou négatif. On dit que
$\mathscr F$ possède la propriété $(S_k)$ si, pour tout $x\in X$, on a
\begin{equation}
\operatorname{prof}(\mathscr F_x)\geq
\inf(k,\dim(\mathscr F_x)).
\tag{5.7.2.1}\label{IV.2.s5.7.2.1-fr}
\end{equation}
On dit que $\mathscr F$ possède la propriété $(S_k)$ en un point $x\in X$
si, pour toute générisation $x'$ de $x$ dans $X$, on a
\begin{equation}
\operatorname{prof}(\mathscr F_{x'})\geq
\inf(k,\dim(\mathscr F_{x'})).
\tag{5.7.2.2}\label{IV.2.s5.7.2.2-fr}
\end{equation}

On dit que $X$ vérifie la propriété $(S_k)$ (resp. vérifie la propriété
$(S_k)$ en un point $x$) si $\mathcal O_X$ vérifie la propriété $(S_k)$
(resp. vérifie la propriété $(S_k)$ au point $x$).

Pour que $\mathscr F$ vérifie la propriété $(S_k)$, il faut et il suffit
évidemment qu'il la vérifie en tout point de $X$. Si $U$ est un ouvert de
$X$ et si $\mathscr F$ vérifie $(S_k)$, il en est de même de
$\mathscr F|U$; réciproquement, si $(U_\alpha)$ est un recouvrement ouvert
de $X$ et si $\mathscr F|U_\alpha$ vérifie $(S_k)$ pour tout $\alpha$,
$\mathscr F$ vérifie $(S_k)$.
\end{definition}

\begin{remarks}[5.7.3]
\label{IV.2.s5.7.3-fr}
\begin{enumerate}
\item[(i)] Rappelons que l'on a toujours
$\operatorname{prof}(\mathscr F_x)\leq\dim(\mathscr F_x)$ (0, 16.4.5.1)
si $\mathscr F_x\neq0$. Dire que $\mathscr F$ possède la propriété $(S_k)$
signifie donc que l'on a $\operatorname{prof}(\mathscr F_x)\geq k$ sauf aux
points $x\in X$ tels que $\dim(\mathscr F_x)<k$ et qu'en ces derniers points
on a $\dim(\mathscr F_x)=\operatorname{prof}(\mathscr F_x)$, c'est-à-dire
(0, 16.5.1) que $\mathscr F$ est un $\mathcal O_X$-Module de
Cohen-Macaulay en ces points; on notera qu'aux points où
$\dim(\mathscr F_x)=k$, on a $\operatorname{prof}(\mathscr F_x)=k$,

